
\section{耿丹学院参考文献格式说明}
\demo{耿丹学院参考文献内容建议：}
\begin{enumerate}
   \item 人文、经管类：中文至少24篇，外文至少6篇；
   \item 工理科类：中文至少20篇，外文至少4篇；
   \item 设计类：中文至少10篇，外文至少2篇。
\end{enumerate}
\begin{demoenv}
参考文献中近五年的文献数一般应不少于总数的1/3，并应有近两年的参考文献。本专业教科书不能作为参考文献。

西文文献中第一个词和每个实词的第一个字母大写，余者小写；俄文文献名第一个词和专有名词的第一个字母大写，余者小写；日文文献中的汉字须用日文汉字，不得用中文汉字、简化汉字代替。文献中的外文字母一律用正体。

作者为多人时，不同作者姓名间用逗号加一空格相隔。外文姓名按国际惯例，将作者名的缩写置前，作者姓置后。

学术刊物文献无卷号的可略去此项，直接写"年，(期)"。

列出作者直接阅读过或在正文中被引用过的文献资料，参考文献要另起一页。

参考文献在文中要有引用标注，引用标注顺序与参考文献先后顺序相同(这个由\LaTeX 自动处理)。

当提及的参考文献为文中直接说明时，则用小四号字与正文排齐，如"由文献[8,10\~14]可知"。

不得将引用文献标示置于论文摘要和各级标题处。

参考文献序号顶格书写，其后空一格写作者名。若同一文献中有多处被引用，则要写出相应引用页码，各起止页码间空一格，没有引用则不写引用起止页。排列按引用顺序，不按页码顺序。
\end{demoenv}
 文献类型和标志代码 
 \begin{table}[htbp]
 \centering
 \caption{文献类型和标志代码}
 \zihao{-5}\songti
 \begin{tabular}{ll}
 \toprule
 文献类型 & 标志代码 \\
 \midrule
 普通图书 & M \\
 会议录 & C \\
 汇编 & G \\
 报纸 & N \\
 期刊 & J \\
 学位论文 & D \\
 报告 & R \\
 标准 & S \\
 专利 & P \\
 数据库 & DB \\
 计算机程序 & CP \\
 电子公告 & EB \\
 \bottomrule
 \end{tabular}
 \end{table}
 
 电子文献载体和标志代码 
 \begin{table}[htbp]
 \centering
 \caption{电子文献载体和标志代码}
 \zihao{-5}\songti
 \begin{tabular}{ll}
 \toprule
 载体类型 & 标志代码 \\
 \midrule
 网络 & OL \\
 磁带 & MT \\
 磁盘 & MK \\
 光盘 & CD \\
 \bottomrule
 \end{tabular}
 \end{table}
